\begin{abstract}
Text-based person search retrieves images of a pedestrian from a large gallery given a free-form description. This report is drawn from an ongoing master's thesis on parameter-efficient fine-tuning (PEFT) for federated text-based person search. In federated learning every client uploads its update in every round, so the size of the trainable part sets the communication cost, and a premise must hold before federating: a far smaller update must lose little accuracy. We test this premise centrally on IRRA, the method of our midterm presentation, which fine-tunes all of CLIP with the Similarity Distribution Matching (SDM) loss. Freezing CLIP and attaching rank-4 LoRA to the attention projections of both encoders, we train \nTrainable{} parameters instead of \nBackboneM\,M, \sizeRatio$\times$ fewer (a \ckptMiB\,MiB instead of a \fullMiB\,MiB update). With SDM as the sole objective the model reaches \sdmPeakR\% Rank-1 on RSTPReid at its best epoch and \sdmConvR\% once converged, on par with the 57.20\% IRRA reports for full fine-tuning with SDM alone; in a \kClients-client federation the per-round upload would fall from \upFullRoundGiB\,GiB to \upLoraRoundMiB\,MiB. We also ask where SDM's gain over the standard contrastive loss comes from. Under the random batch sampling IRRA uses, only \posFrac\% of batch rows contain a second positive, so SDM's target is almost always one-hot; there SDM does \emph{not} reduce to InfoNCE but to a scaled mean-absolute-error term plus an entropy bonus. A control loss with SDM's targets and the opposite KL direction tracks InfoNCE to within \mpNceDiffR{} Rank-1 points on average and trails SDM by \gapSdmMpceR{} points, attributing the gain to the divergence direction rather than to multi-positive supervision.
\end{abstract}

\begin{IEEEkeywords}
text-based person search, cross-modal retrieval, CLIP, LoRA, parameter-efficient fine-tuning, federated learning, contrastive loss
\end{IEEEkeywords}
